\subsubsection{E-handel}
\label{sec:ehandel}
E-handel er, i modsætning til detailhandel, elektronisk handel via internettet \cite{ddo_ehandel}. På internettet kan lampebutikker have såkaldte e-butikker, hvor kunder kan købe varer \cite{ddo_ebutik}. E-butikker er ofte udformet således, at kunden kan se billeder og informationer omkring lampebutikkens varer og derudfra kan kunden vælge at lægge varerne i en virtuel indkøbskurv, hvor kunden til sidst indtaster de nødvendige oplysninger for at købe og modtage varerne.

Blandt de mange forskellige varer, der sælges via e-butikker, er det her relevant at tale om e-handel med lamper. Nedenstående figur \ref{fig:e_handel_med_lamper} illustrerer princippet bag en lampesælgers salg af lampe til en kunde via en e-butik.

\begin{figure}[H]
  \includegraphics{e_handel_med_lampe.pdf}
  \caption{Princippet bag handel af en lampe via en e-butik.}
    \label{fig:e_handel_med_lamper}
\end{figure}

På figur \ref{fig:e_handel_med_lamper} er det vist, hvordan e-handlen starter med, at kunden får et udvalg af lamper fra e-butikken. Kunden sender så en bestilling, som via e-butikken sendes videre til lampebutikken, og til sidst sendes lampen til kunden. Dog ender handlen ikke nødvendigvis her, da kunden kan sende lampen retur såfremt at de gældende lovgivninger og købsbetingelser muliggør dette. For at undersøge lovgivningen nærmere kan man tage udgangspunkt i den danske lov om forbrugeraftaler, nærmere bestemt ved: LOV nr 1457 af 17/12/2013 Gældende (Forbrugeraftaleloven), Offentliggørelsesdato: 18-12-2013
Justitsministeriet. \cite{retsinformationen}.

I lovens kapitel 1, § 1, stk.\ 2, nr.\ 1 fremgår det, at lovens bestemmelser for fortrydelsesret gælder for aftaler, som er indgået ved fjernsalg. For en  fjernsalgsaftale gælder der, at aftalen om varer, er indgået gennem fjernkommunikation, hvor den erhvervsdrivende og forbrugeren ikke mødes fysisk (jf.\ kap. 1, § 3, nr.\ 1).

Ser man nu på loven i forbindelse med e-handel, foregår fjernkommunikationen gennem internettet via e-butikken, hvor fjernsalgsaftalen udføres i form af brugerens bestilling af f.eks.\ en lampe. Dette gør at fortrydelsesretten gælder ved e-handel.

Fortrydelsesretten er en kundes mulighed for at melde sig ud af en aftale, herunder køb af lamper ved e-handel. Hvis en kunde eksempelvis køber en lampe via en e-butik, har kunden mulighed for at fortryde købet inden 14 dage ved at meddele dette til den erhvervsdrivende (jf.\ kap.\ 4, § 19). Herefter har kunden 14 dage til at returnere varen (jf.\ kap.\ 4, § 24). Hvis varens værdi er forringet som følge af kundens unødvendige håndtering af varen for at inspicere denne, så hæfter kunden for værdiforringelsen (jf.\ kap.\ 4, § 24, stk.\ 5). Dvs.\ at hvis en bruger installerer og bruger lampen, hvor der f.eks.\ tilpasses ledninger, så kan lampens værdi forringes og kunden skal hæfte for dette. 
